
\documentclass[UTF8,11pt,a4paper]{ctexart}
\usepackage[a4paper,left=23mm,right=23mm,top=24mm,bottom=24mm]{geometry}
\usepackage{fontspec}
\setmainfont{DejaVu Sans}
\setCJKmainfont{Noto Serif CJK SC}
\usepackage{amsmath,amssymb}
\usepackage{xcolor}
\usepackage[colorlinks=true,linkcolor=blue!45!black,unicode=true]{hyperref}
\usepackage{enumitem}
\usepackage{needspace}
\usepackage{fancyhdr}
\setlength{\parindent}{2em}
\setlength{\parskip}{0.25em}
\setlist[enumerate]{leftmargin=2.6em,itemsep=0.4em,topsep=0.5em}
\pagestyle{fancy}
\fancyhf{}
\fancyhead[L]{\small 三千弱水 · 大学化学高阶综合题库}
\fancyhead[R]{\small 待专家审校稿}
\fancyfoot[C]{\thepage}
\setlength{\headheight}{15pt}
\sloppy

\begin{document}
\begin{center}
{\LARGE\bfseries 大学化学高阶综合题 · 第11辑 (12题)}\\[0.8em]
{\large 题目、分步解析与数值核验}\\[0.5em]
{\small \textcolor{red!65!black}{审校稿：未经化学专家核准；禁止直接视为正式试卷或出版物}}
\end{center}
\vspace{1em}
\noindent 本册从仓库结构化 JSONL 自动生成。增补和修改应只编辑源题目数据，并重新执行统一编译。
\par 机械方程残差检查只能说明指定模型的内部数值一致，不保证模型现实可行性。第六辑另须完成跨辑模型去重审查。
\tableofcontents
\clearpage
\section*{第11辑 · 题目及解答}
\addcontentsline{toc}{section}{第11辑}
\subsection*{无机化学}
\addcontentsline{toc}{subsection}{第11辑 / 无机化学}
\Needspace{7\baselineskip}
\subsubsection*{UC-ADV11-I126 · SrFeO₃₋δ氧缺陷的TGA反演、铁价态守恒与晶胞密度}
\noindent\textbf{题目}\par
假想钙钛矿SrFeO₃₋δ从完全氧化的SrFeO₃（Fe全为+4）在还原气氛下失去氧，热重分析显示相对初始样品质量损失恰为1.000\%。取原子量Sr=87.620、Fe=55.845、O=15.999 g/mol；晶胞为立方钙钛矿、每晶胞1个化学式单位、a=0.3900 nm，NA=6.02214076×10²³ mol⁻¹。假定Sr始终+2、O为−2、Fe只存在+3/+4，所有氧损失形成氧空位且无阳离子缺位或二次相。求氧缺陷δ、Fe³⁺与Fe⁴⁺的位点分数、平均铁价、氧空位/晶胞、还原后式量与理论密度；用电中性和热重质量差反向核对，并说明可能的物理模型失效。\par
\medskip\noindent\textbf{分步解答}
\begin{enumerate}
\item 初始SrFeO₃式量M₀=87.620+55.845+3×15.999=191.462000000 g/mol；质量损失必须相对于初始式量计算。
\item 每化学式单位失去δ个氧，损失质量δ×15.999 g/mol，故δ=0.01000M₀/15.999=0.119671229452。
\item 设Fe⁴⁺位点分数y，平均铁价为3+y；电中性要求2+(3+y)−2(3−δ)=0，因此y=1−2δ=0.760657541096。
\item Fe³⁺分数为1−y=0.239342458904；平均铁价为3.76065754110，严格处于3与4之间。
\item 每个晶胞1个式量单位，氧空位数/晶胞=δ=0.119671229452；晶胞中氧平均占位数为3−δ。
\item 还原后摩尔质量M=M₀−15.999δ=189.547380000 g/mol；与初始差值占M₀的比例严格为1.000\%。
\item 晶胞体积a³=(0.3900×10⁻⁹ m)³，理论密度ρ=(M/1000)/(NAa³)=5306.07106611 kg/m³，必须使用SI单位。
\item 氧空位反应每去除一个O²⁻对应两个电子，恰将2δ份Fe⁴⁺还原为Fe³⁺，与电中性方程一致。
\item 若δ>0.5会要求Fe⁴⁺负分数，违反本题价态假设；真实钙钛矿可能发生结构相变、Fe²⁺、空位有序或多相析出，需要独立衍射与谱学证据。
\end{enumerate}
\noindent\textbf{答案摘要}\par
delta = 0.119671229 per\_formula；Fe4\_fraction = 0.760657541 fraction；Fe3\_fraction = 0.239342459 fraction；Fe\_mean\_valence = 3.76065754 oxidation\_state；vacancies\_per\_cell = 0.119671229 per\_cell；formula\_mass = 189.54738 g/mol；density = 5306.07107 kg/m³\par
\noindent\textbf{模型前提：}
单相立方SrFeO₃₋δ、Sr2+、Fe3+/4+、氧空位独占；实际晶格常数随δ变化被忽略。\par
\noindent\textcolor{red!65!black}{[跨辑去重未完成 · 草稿]} \par
\medskip\hrule\medskip
\Needspace{7\baselineskip}
\subsubsection*{UC-ADV11-I127 · 二价钙置换沸石钠位点的双位点质量作用与van’t Hoff选择性}
\noindent\textbf{题目}\par
假想沸石有0.01000 mol负电荷位点，Na⁺每个占1位，Ca²⁺每个占2位。与无限大理想储液池平衡，aNa=0.1000、aCa=0.02000均为相对1 mol/L的无量纲活度。设θNa为被Na占据的位点分数、θCa=1−θNa为被Ca占据的位点分数；定义理想化交换常数K=θCa aNa²/(θNa² aCa)。298.15 K时K=5.000，交换反应标准焓ΔH°=−20.00 kJ/mol，假定298–318 K恒定。求298.15 K的两类位点分数、沸石中Na/Ca物质的量、Ca占位分数、318.15 K的K与分数；检查电荷/位点守恒、正根唯一性及放热交换的温度趋势。\par
\medskip\noindent\textbf{分步解答}
\begin{enumerate}
\item 位点守恒θNa+θCa=1，Ca每离子占两负位，因此nNa=0.01000θNa，nCa=0.01000θCa/2。
\item 由定义K=θCa aNa²/(θNa² aCa)整理为(1−θNa)/θNa²=K aCa/aNa²=q。
\item 298.15 K时q=10.00000000，方程qθNa²+θNa−1=0，唯一正根θNa=2/[1+√(1+4q)]=0.270156211872。
\item θCa=1−θNa=0.729843788128；负根不在[0,1]，故排除。
\item 沸石Na=0.00270156211872 mol，Ca=0.00364921894064 mol；电荷当量nNa+2nCa=0.01000 mol。
\item van’t Hoff积分ln(K₂/K₁)=(ΔH°/R)(1/T₁−1/T₂)，升温至318.15 K得K₂=3.01096757128，放热交换升温使K减小。
\item 新q₂=K₂aCa/aNa²=6.02193514255，对应θNa₂=0.332847023191、θCa₂=0.667152976809。
\item 方程左边(1−θ)/θ²对0<θ≤1严格递减，因此每个正q仅有一个物理解；温升时Ca选择性降低。
\item 该选择性常数使用题设理想位点分数而非真实固相活度，未计入邻位关联、交换热变化、溶液非理想和位点能分布。
\end{enumerate}
\noindent\textbf{答案摘要}\par
q298 = 10 dimensionless；theta\_Na298 = 0.270156212 fraction；theta\_Ca298 = 0.729843788 fraction；nNa298 = 0.00270156212 mol；nCa298 = 0.00364921894 mol；K318 = 3.01096757 dimensionless；q318 = 6.02193514 dimensionless；theta\_Na318 = 0.332847023 fraction；theta\_Ca318 = 0.667152977 fraction\par
\noindent\textbf{模型前提：}
二价离子占两位、固定外液活度、理想位点混合、ΔH°恒定；非理想沸石选择性仍待专家复核。\par
\noindent\textcolor{red!65!black}{[跨辑去重未完成 · 草稿]} \par
\medskip\hrule\medskip
\Needspace{7\baselineskip}
\subsubsection*{UC-ADV11-I128 · Cu(II)轴向Jahn–Teller畸变的EPR各向异性与晶场跃迁能反演}
\noindent\textbf{题目}\par
一假想轴向伸长Cu(II) d⁹配合物的低温EPR主轴g因子为g∥=2.3000、g⊥=2.0700，自由电子ge=2.0023。按指定的一阶自旋轨道微扰模型，g∥=ge−8λ/Δ₁、g⊥=ge−2λ/Δ₂，λ=−830.0 cm⁻¹，Δ₁、Δ₂为相应激发态能隙（单位cm⁻¹）；假定微扰适用且忽略共价性折减、轨道混合与超精细项。求两个晶场能隙及相应波长λ光(nm)=10⁷/Δ(cm⁻¹)、粉末取向平均ḡ、g的均方根、S=1/2有效磁矩μeff=g\_rms√[S(S+1)]μB、g各向异性；核对能隙正值和微扰展开参数，并说明实验局限。\par
\medskip\noindent\textbf{分步解答}
\begin{enumerate}
\item 由g∥−ge=−8λ/Δ₁反解Δ₁=−8λ/(g∥−ge)=22304.3332214 cm⁻¹；g∥>ge与负λ使Δ₁>0。
\item 由g⊥−ge=−2λ/Δ₂得到Δ₂=24519.9409158 cm⁻¹；不可将两条微扰式的系数8与2互换。
\item 波数转换为真空波长，λ₁=10⁷/Δ₁=448.343373494 nm，λ₂=407.831325301 nm。
\item 随机粉末的简单算术取向平均ḡ=(g∥+2g⊥)/3=2.14666666667，但磁化率应采用平方平均而非ḡ²。
\item 平方平均g\_rms=√[(g∥²+2g⊥²)/3]=2.14940301789，对应S=1/2的μeff=1.86143761647 μB。
\item g各向异性Δg=g∥−g⊥=0.230000000000，轴向伸长的特征需结合配体场与真实轨道基态判别。
\item 微扰量级|λ|/Δ₁=0.03721250000、|λ|/Δ₂=0.03385000000，均显著小于1，支持指定一阶近似的内部自洽。
\item 回代g∥=ge−8λ/Δ₁与g⊥=ge−2λ/Δ₂严格复现给定g值；本反演只在假定轨道模型下唯一。
\item 实际Cu(II)的共价性折减、超精细、低对称场混合及动态Jahn–Teller会改变系数，不能直接把两个波数当作已实测的吸收峰。
\end{enumerate}
\noindent\textbf{答案摘要}\par
Delta1 = 22304.3332 cm\textasciicircum{}-1；Delta2 = 24519.9409 cm\textasciicircum{}-1；wavelength1 = 448.343373 nm；wavelength2 = 407.831325 nm；g\_average = 2.14666667 relative；g\_rms = 2.14940302 relative；mu\_eff = 1.86143762 Bohr\_magneton；g\_anisotropy = 0.23 relative；lambda\_over\_Delta1 = 0.0372125 dimensionless\par
\noindent\textbf{模型前提：}
给定轴向d9自旋轨道微扰模型，忽略共价性与低对称修正；与历史EPR模型需人工逐式比对。\par
\noindent\textcolor{red!65!black}{[跨辑去重未完成 · 草稿]} \par
\medskip\hrule\medskip
\subsection*{分析化学}
\addcontentsline{toc}{subsection}{第11辑 / 分析化学}
\Needspace{7\baselineskip}
\subsubsection*{UC-ADV11-A126 · 电化学石英晶体微天平的铜沉积法拉第质量与水合数反演}
\noindent\textbf{题目}\par
一块5 MHz石英晶体电极在水溶液中还原Cu²⁺+2e⁻→Cu(s)，电极有效面积1.000 cm²。沉积过程通过电量Q=0.001000 C，晶体基频变化Δf=−34.00 Hz；给定Sauerbrey质量灵敏度Cf=56.60 Hz/(μg/cm²)，法拉第常数F=96485.33212 C/mol，Cu摩尔质量63.546 g/mol，水18.015 g/mol。假设沉积法拉第效率100\%、膜刚性且均匀、无其他吸附离子，频移只由Cu和共沉积水贡献。求总面质量、Cu物质的量及质量、附带水质量及物质的量、每Cu对应的水分子数、Cu面摩尔密度，并回算频移和电子守恒，说明何时Sauerbrey反演失效。\par
\medskip\noindent\textbf{分步解答}
\begin{enumerate}
\item Sauerbrey关系Δf=−Cf Δm/A，故总面质量=34.00/56.60=0.600706713781 μg/cm²。
\item 每个Cu²⁺消耗2电子，nCu=Q/(2F)=5.18213482831e-9 mol，不能误用单电子反应。
\item 电极面积为1.000 cm²，Cu面摩尔密度nCu/A=5.18213482831e-9 mol/cm²。
\item 法拉第理论Cu质量mCu=nCu×63.546 g/mol=0.329303939800 μg，低于QCM总面质量。
\item 差额归于共沉积水，mwater=mtotal−mCu=0.271402773981 μg；正值支持题设水合解释，但不是其独立实验验证。
\item 水物质的量nwater=mwater×10⁻⁶/18.015=1.50653774067e-8 mol。
\item 每个Cu对应水合数h=nwater/nCu=2.90717588519，可为非整数平均值。
\item 将mCu+mwater代回Δf=−Cf m/A严格恢复−34.00 Hz；电子总量2FnCu=0.001000 C。
\item 若膜黏弹性显著、液体黏度改变、存在夹带阴离子或析氢，频移不能仅解释为刚性铜水合膜，需QCM-D耗散与法拉第效率独立测定。
\end{enumerate}
\noindent\textbf{答案摘要}\par
total\_ug\_per\_cm2 = 0.600706714 μg/cm²；nCu = 5.18213483e-09 mol；mCu\_ug = 0.32930394 μg；mWater\_ug = 0.271402774 μg；nWater = 1.50653774e-08 mol；hydration\_number = 2.90717589 water/Cu；Cu\_mol\_per\_cm2 = 5.18213483e-09 mol/cm²；frequency\_back = -34 Hz\par
\noindent\textbf{模型前提：}
刚性膜、无溶液黏度变化、100\%铜法拉第效率且只夹带水；真实QCM-D需耗散验证。\par
\noindent\textcolor{red!65!black}{[跨辑去重未完成 · 草稿]} \par
\medskip\hrule\medskip
\Needspace{7\baselineskip}
\subsubsection*{UC-ADV11-A127 · 高分辨质谱¹³C与¹⁵N精细结构分辨及单同位素峰强度推断}
\noindent\textbf{题目}\par
一分子含C₂₀N₂且不含其他可贡献M+1的同位素，基峰质荷比m/z=500.0000（z=+1）。相对¹²C的¹³C质量增量为1.003354835 u，相对¹⁴N的¹⁵N增量为0.997034893 u；天然丰度p(¹³C)=0.01070、p(¹⁵N)=0.003640，假设各原子同位素独立。质谱在m/z 500处分辨本领R=m/ΔmFWHM=120000，且各峰FWHM相同。计算单¹³C与单¹⁵N取代峰的精细质量差、FWHM、峰距/FWHM、最低使峰距等于FWHM的R、两个单取代同位素峰的概率比；若离子为+2且m/z约250、分辨本领相同，求峰距与FWHM并讨论为何质量差与分辨率不能混用。\par
\medskip\noindent\textbf{分步解答}
\begin{enumerate}
\item 单¹³C取代比全轻同位素峰高1.003354835 u，单¹⁵N取代高0.997034893 u，二者精细质量差=0.00631994200000 u。
\item 单电荷离子Δ(m/z)等于质量差；仪器FWHM=500/120000=0.00416666666667 u。
\item 峰距与峰宽之比=1.51678608000；数值大于1仅说明峰心相距超过一个FWHM，不能保证基线分离。
\item 临界Rmin=500/Δm=79114.6501028，须区分仪器标称分辨本领与实际峰形。
\item 恰好一个¹³C且无¹⁵N的概率P13=20p13(1−p13)\textasciicircum{}19(1−p15)\textasciicircum{}2，因20个碳位可被取代。
\item 恰好一个¹⁵N且无¹³C的概率P15=2p15(1−p15)(1−p13)\textasciicircum{}20；两者概率比P13/P15=29.6053819828。
\item 若电荷为+2，精细峰距为Δm/2=0.00315997100000 m/z，而m/z约250处同R的FWHM=250/120000=0.00208333333333，两者比值保持不变。
\item 精细同位素峰的相对丰度与质量分辨是两类独立信息；共洗脱杂质、质量偏差、其它元素同位素及峰尾会改变实际识别。
\item 概率计算是指定C₂₀N₂模型的条件预测，不能仅凭一个质量差断言分子式；仍需母离子准确质量、元素约束和同位素包络整体拟合。
\end{enumerate}
\noindent\textbf{答案摘要}\par
fine\_mass\_gap = 0.006319942 u；FWHM\_1plus = 0.00416666667 u；gap\_over\_FWHM = 1.51678608 ratio；R\_min = 79114.6501 ratio；P13\_over\_P15 = 29.605382 ratio；gap\_2plus = 0.003159971 m/z；FWHM\_2plus = 0.00208333333 m/z\par
\noindent\textbf{模型前提：}
只考虑碳氮两种重同位素、独立天然丰度、相同峰形与恒定分辨本领；非真实质谱鉴定结论。\par
\noindent\textcolor{red!65!black}{[跨辑去重未完成 · 草稿]} \par
\medskip\hrule\medskip
\Needspace{7\baselineskip}
\subsubsection*{UC-ADV11-A128 · 时间分辨FRET寿命反演的供受体距离与双受体竞争转移}
\noindent\textbf{题目}\par
一假想标记蛋白的荧光供体在无受体时寿命τD=4.000 ns、荧光量子产率ΦD=0.8000；引入一个固定受体后寿命τDA=1.600 ns。假定受体只引入FRET非辐射转移通道，供体自身辐射与其他非辐射速率均不变，取Förster半径R₀=5.000 nm、取向因子已含在R₀中。单受体效率E=1−τDA/τD=1/[1+(r/R₀)⁶]。求转移效率、转移速率、供受体距离、单受体条件下供体量子产率；若两个相同受体各自距供体r且独立竞争转移，求总寿命、总FRET效率、量子产率并核查速率相加与适用条件。\par
\medskip\noindent\textbf{分步解答}
\begin{enumerate}
\item 无受体供体总失活速率kD=1/τD=0.2500 ns⁻¹，其中辐射与固有非辐射通道均包含在内。
\item 单受体总失活速率kDA=1/τDA=0.6250 ns⁻¹，新增FRET速率kET=kDA−kD=0.375000000000 ns⁻¹。
\item 转移效率E=kET/(kD+kET)=1−τDA/τD=0.600000000000，严格位于0与1之间。
\item Förster距离律E=1/[1+(r/R₀)⁶]，反演r=R₀(1/E−1)\textasciicircum{}(1/6)=4.67327632592 nm。
\item 若供体辐射速率不变，量子产率与寿命成正比，ΦDA=ΦD×τDA/τD=0.320000000000。
\item 两个相同受体独立提供两个kET通道，总失活速率k2=kD+2kET，故τ₂=1.00000000000 ns。
\item 双受体总效率E₂=2kET/(kD+2kET)=0.750000000000，不能简单将单受体效率乘2而超过1。
\item 双受体供体量子产率Φ₂=ΦD×τ₂/τD=0.200000000000，寿命和量子产率均减小。
\item 本模型要求单指数寿命、固定距离和取向、无直接受体激发及能量回传；蛋白构象分布或光谱串色会使反演距离不唯一。
\end{enumerate}
\noindent\textbf{答案摘要}\par
donor\_decay\_rate = 0.25 ns\textasciicircum{}-1；FRET\_rate = 0.375 ns\textasciicircum{}-1；FRET\_efficiency = 0.6 fraction；distance = 4.67327633 nm；quantum\_yield\_one = 0.32 fraction；lifetime\_two = 1 ns；FRET\_eff\_two = 0.75 fraction；quantum\_yield\_two = 0.2 fraction\par
\noindent\textbf{模型前提：}
固定取向和距离、独立转移通道、辐射速率不变；真实寿命分布与光谱串色待实验审核。\par
\noindent\textcolor{red!65!black}{[跨辑去重未完成 · 草稿]} \par
\medskip\hrule\medskip
\subsection*{有机化学}
\addcontentsline{toc}{subsection}{第11辑 / 有机化学}
\Needspace{7\baselineskip}
\subsubsection*{UC-ADV11-O123 · 活性阴离子聚合的Poisson链长分布、转化率与分散系数闭合}
\noindent\textbf{题目}\par
假想活性阴离子聚合体系中单体初始浓度[M]₀=1.200 mol/L，瞬时且定量生成活性链端[I*]=0.01200 mol/L，每条活性链端独立以相同速率增长、无终止和链转移，增长常数kp=0.05000 L/(mol s)。定义单体转化率X=0.8000，忽略体积变化，重复单元摩尔质量104.0 g/mol。将链中增长单元数n近似为参数λ=消耗单体浓度/[I*]的Poisson分布（可有概率极小的n=0），忽略引发端基质量。求达到80\%转化率的时间、剩余单体、平均链长、数均与重均分子量、分散系数Đ、活性链端浓度和链长方差，解释为何活性链端数不随单体消耗下降，并验证Poisson二阶矩。\par
\medskip\noindent\textbf{分步解答}
\begin{enumerate}
\item 无终止和链转移时活性链端浓度[I*]恒定为0.01200 mol/L，不能当作单体浓度一起衰减。
\item 单体消耗动力学d[M]/dt=−kp[I*][M]，积分得[M](t)=M₀exp(−kp[I*]t)。
\item 转化率0.8意味着[M]=1.200×0.2000=0.240000000000 mol/L。
\item 达到该转化率的时间t=−ln(0.2)/(0.05000×0.01200)=2682.39652072 s。
\item 消耗单体浓度M₀X=0.9600 mol/L，平均每条链增长单元数λ=0.9600/0.01200=80.0000000000。
\item Poisson分布E[n]=λ、Var(n)=λ、E[n²]=λ²+λ；因此数均聚合度DPn=λ，重均聚合度DPw=E[n²]/E[n]=λ+1。
\item 数均分子量Mn=104λ=8320.00000000 g/mol，重均Mw=104(λ+1)=8424.00000000 g/mol。
\item 分散系数Đ=Mw/Mn=1+1/λ=1.01250000000，接近1但不等于1。
\item 由于n=0的概率exp(−80)极小，Poisson近似在本题高转化下合理；真实引发时间分布、离子对聚集、链转移与端基质量会改变链长分布。
\end{enumerate}
\noindent\textbf{答案摘要}\par
remaining\_monomer = 0.24 mol/L；time\_to\_80 = 2682.39652 s；lambda = 80 units；variance\_n = 80 units²；DPw = 81 units；Mn = 8320 g/mol；Mw = 8424 g/mol；dispersity = 1.0125 ratio；active\_chain\_ends = 0.012 mol/L\par
\noindent\textbf{模型前提：}
同步定量引发、无终止链转移、Poisson独立增长、忽略端基质量；与v10自由基终止动力学是不同核心随机模型，仍待人工跨辑核验。\par
\noindent\textcolor{red!65!black}{[跨辑去重未完成 · 草稿]} \par
\medskip\hrule\medskip
\Needspace{7\baselineskip}
\subsubsection*{UC-ADV11-O124 · 二氯一溴有机物同位素包络的生成函数与重同位素矩检验}
\noindent\textbf{题目}\par
某假想有机物分子含恰好2个氯和1个溴，其他元素同位素贡献可忽略。设³⁷Cl天然丰度pCl=0.2422、³⁵Cl丰度0.7578，⁸¹Br丰度pBr=0.4931、⁷⁹Br丰度0.5069，各原子独立且忽略同位素质量缺陷，将每次重同位素替换视为质量约增加2 u。求M、M+2、M+4、M+6四个峰的归一化概率、M+6与M峰强度比、至少出现一个重同位素的概率、重同位素数的均值；使用概率生成函数和导数核验，解释为何不能把Cl和Br视作同一丰度的三个二项试验。\par
\medskip\noindent\textbf{分步解答}
\begin{enumerate}
\item 对每个氯原子，生成多项式qCl+pCl z，其中qCl=0.7578；两个独立氯给平方。
\item 一个溴原子对应qBr+pBr z，qBr=0.5069；总生成函数G(z)=(qCl+pCl z)²(qBr+pBr z)。
\item 常数项P(M)=qCl²qBr=0.291092819796。
\item z的一次项P(M+2)=2pClqClqBr+qCl²pBr=0.469240020612。
\item 二次项P(M+4)=pCl²qBr+2pClqClpBr=0.210741499388。
\item 三次项P(M+6)=pCl²pBr=0.0289256602040；四个系数和G(1)=1。
\item 高质量端与单同位素端强度比P(M+6)/P(M)=0.0993691985404，至少一重同位素概率=1−P(M)=0.708907180204。
\item 由G′(1)=2pCl+pBr得到重同位素数期望=0.977500000000，与四个峰的k加权和一致。
\item 氯与溴重同位素丰度不同，不能用单一Binomial(3,p)；实际高分辨质谱需加入精确质量、其他元素同位素与仪器响应。
\end{enumerate}
\noindent\textbf{答案摘要}\par
P\_M = 0.29109282 fraction；P\_M2 = 0.469240021 fraction；P\_M4 = 0.210741499 fraction；P\_M6 = 0.0289256602 fraction；P\_M6\_over\_M = 0.0993691985 ratio；P\_at\_least\_one = 0.70890718 fraction；mean\_heavy = 0.9775 count\par
\noindent\textbf{模型前提：}
独立天然同位素、Cl2Br计数已知、质量增量仅按2u近似；真实质谱需精确质量修正。\par
\noindent\textcolor{red!65!black}{[跨辑去重未完成 · 草稿]} \par
\medskip\hrule\medskip
\Needspace{7\baselineskip}
\subsubsection*{UC-ADV11-O125 · Fischer酯化平衡与中途除水的二次反应进度及物料闭合}
\noindent\textbf{题目}\par
在298.15 K的理想液相中，假想乙酸A与乙醇B发生A+B⇌乙酸乙酯E+水W，平衡常数K=(nE nW)/(nA nB)=4.000，假定各物种理想且总体积不变。初始nA=1.000 mol、nB=2.000 mol、nE=0.2000 mol、nW=0.1000 mol。先在密闭条件下达到平衡；随后在不改变A/B/E物质的量的条件下瞬间移除恰好一半的平衡水，再重新密闭至同温平衡。求首次反应进度、四物种库存、除水量、第二次追加反应进度、最终酯量和酸转化率；证明两次物理根唯一性、平衡移动方向及乙酸/乙醇反应计量守恒。\par
\medskip\noindent\textbf{分步解答}
\begin{enumerate}
\item 令首次正向进度ξ，库存nA=1−ξ、nB=2−ξ、nE=0.2+ξ、nW=0.1+ξ，物理区间0≤ξ<1。
\item 首次反应商Q(ξ)=(0.2+ξ)(0.1+ξ)/[(1−ξ)(2−ξ)]；Q(0)=0.0100<K，故向右反应。
\item 解Q(ξ)=4得首次进度ξ₁=0.808025765163 mol；Q随ξ严格递增，故物理解唯一。
\item 首次平衡库存A=0.191974234837、B=1.19197423484、E=1.00802576516、W=0.908025765163 mol。
\item 瞬间移除一半平衡水，移除量=0.454012882581 mol，剩余水也为0.454012882581 mol；A、B、E在移除瞬间不变。
\item 除水后Q变为原来的一半，即2<K，故体系继续向右酯化，设第二次追加进度η≥0。
\item 第二次平衡式为[(0.2+ξ₁+η)(0.5(0.1+ξ₁)+η)]/[(1−ξ₁−η)(2−ξ₁−η)]=4，解η=0.0673440443911 mol。
\item 最终总反应进度ξtot=ξ₁+η=0.875369809554 mol，酯E=1.07536980955 mol，初始酸转化率=0.875369809554。
\item 酸剩余=0.124630190446 mol，乙醇剩余=1.12463019045 mol；除水操作不改变A、B、E计量守恒关系nA+nE=1.2、nB+nE=2.2。真实非理想活度及体积变化需修正K表达式。
\end{enumerate}
\noindent\textbf{答案摘要}\par
xi\_first = 0.808025765 mol；acid\_first = 0.191974235 mol；alcohol\_first = 1.19197423 mol；ester\_first = 1.00802577 mol；water\_first = 0.908025765 mol；water\_removed = 0.454012883 mol；xi\_second = 0.0673440444 mol；xi\_total = 0.87536981 mol；ester\_final = 1.07536981 mol；acid\_conversion = 0.87536981 fraction；acid\_final = 0.12463019 mol\par
\noindent\textbf{模型前提：}
理想液相、等温固定K、除水瞬间不带走有机组分；真实酯化需考虑活度及酸催化动力学。\par
\noindent\textcolor{red!65!black}{[跨辑去重未完成 · 草稿]} \par
\medskip\hrule\medskip
\subsection*{物理化学}
\addcontentsline{toc}{subsection}{第11辑 / 物理化学}
\Needspace{7\baselineskip}
\subsubsection*{UC-ADV11-P123 · Semenov放热反应器双稳态根与热失控临界参数}
\noindent\textbf{题目}\par
某理想连续补给放热反应器采用Semenov集总温升模型：令θ为无量纲温升，热生成功率Qgen=Q₀exp(θ)，散热功率Qloss=Hθ，H=40.00 W，Q₀=10.00 W。忽略物料耗尽、反应物浓度变化和热容温度依赖，热容C>0满足C dθ/dt=Qgen−Qloss。定义δ=Q₀/H。求全部非负稳态温升、各稳态功率、线性稳定性、临界δ与临界Q₀，并证明当Q₀超过临界值时无有限非负稳态。使用二分法或Lambert W分支验证两根，说明模型适用限制。\par
\medskip\noindent\textbf{分步解答}
\begin{enumerate}
\item 稳态Q₀expθ=Hθ，整理为θexp(−θ)=δ=10/40=0.2500000000。
\item 函数f(θ)=θexp(−θ)在θ≥0先增后减，f′=(1−θ)exp(−θ)，最大值f(1)=1/e。
\item 因为δ=0.25<1/e，存在两个非负根，分别位于(0,1)与(1,∞)，而非只有一个解。
\item 对(0,1)区间二分求低温根θlow=0.357402956181；对(1,10)二分求高温根θhigh=2.15329236411。
\item 两个稳态功率均满足Qgen=Qloss=Hθ，低根功率=14.2961182473 W，高根功率=86.1316945644 W。
\item 线性化热收支导数Q₀expθ−H，在稳态等于H(θ−1)；低根θ<1导数为负、线性稳定，高根θ>1导数为正、线性不稳定。
\item 鞍结临界点θc=1、δc=1/e=0.367879441171，临界Q₀,c=H/e=14.7151776469 W。
\item 当Q₀>H/e时δ>1/e，θexp(−θ)≤1/e无解；模型预测无有限稳态并非证明真实反应器温度无界。
\item 模型忽略反应物耗尽、热传递系数温变与相变，不能将高温不稳定根误称第二个稳定运行状态。
\end{enumerate}
\noindent\textbf{答案摘要}\par
delta = 0.25 dimensionless；theta\_low = 0.357402956 dimensionless；theta\_high = 2.15329236 dimensionless；Q\_low = 14.2961182 W；Q\_high = 86.1316946 W；delta\_critical = 0.367879441 dimensionless；Q0\_critical = 14.7151776 W；theta\_critical = 1 dimensionless\par
\noindent\textbf{模型前提：}
集总等温场、固定反应物、指数热生成与线性散热；真实热失控受物料与热容反馈约束。\par
\noindent\textcolor{red!65!black}{[跨辑去重未完成 · 草稿]} \par
\medskip\hrule\medskip
\Needspace{7\baselineskip}
\subsubsection*{UC-ADV11-P124 · 单分子结合两态Markov过程的停留时间、相关函数与谱密度}
\noindent\textbf{题目}\par
单分子受体在恒定配体浓度c=2.000 μmol/L下发生U⇌B，结合速率kon=2.000×10⁶ L/(mol s)，解离速率koff=6.000 s⁻¹；假设配体池无限大、受体无协同作用、两态连续时间Markov过程。定义结合态指示量X(t)∈\{0,1\}。求结合/解离跃迁率、平衡结合概率、Kd、平均结合与未结合停留时间、弛豫时间τ、X的方差、延迟0.2000 s的平衡自协方差、单边角频率功率谱S(ω)=2Var(X)τ/[1+(ωτ)²]在ω=5.000 rad/s的值；用详细平衡和概率归一性核验。\par
\medskip\noindent\textbf{分步解答}
\begin{enumerate}
\item 将c=2.000×10⁻⁶ mol/L代入，未结合到结合的伪一级速率a=kon c=4.00000000000 s⁻¹。
\item 结合到未结合速率b=koff=6.00000000000 s⁻¹；主方程dpB/dt=a(1−pB)−bpB。
\item 稳态pB=a/(a+b)=0.400000000000，pU=b/(a+b)=0.600000000000，两者之和为1。
\item 解离常数Kd=koff/kon=0.00000300000000000 mol/L；pB=c/(Kd+c)与主方程一致。
\item 两态停留时间均为指数分布，平均结合停留1/b=0.166666666667 s，平均未结合停留1/a=0.250000000000 s。
\item 非零动力学特征值为−(a+b)，故相关弛豫时间τ=1/(a+b)=0.100000000000 s。
\item 二值变量方差pB(1−pB)=0.240000000000；延迟t=0.2 s的协方差C(t)=Var exp(−t/τ)=0.0324804679768。
\item 在ω=5 rad/s，单边角频率谱S(ω)=2Varτ/(1+ω²τ²)=0.0384000000000 s，单位随角频率谱约定而定。
\item 详细平衡pU a=pB b=2.40000000000 s⁻¹；若存在多配体、构象转换或非Markov停留时间，单指数相关与洛伦兹谱失效。
\end{enumerate}
\noindent\textbf{答案摘要}\par
a = 4 s\textasciicircum{}-1；b = 6 s\textasciicircum{}-1；p\_bound = 0.4 fraction；p\_unbound = 0.6 fraction；Kd = 3e-06 mol/L；bound\_dwell = 0.166666667 s；unbound\_dwell = 0.25 s；tau = 0.1 s；variance = 0.24 dimensionless；covariance\_0p2 = 0.032480468 dimensionless；spectrum\_omega5 = 0.0384 s；detailed\_balance\_flux = 2.4 s\textasciicircum{}-1\par
\noindent\textbf{模型前提：}
两态连续时间Markov、恒定外液、无协同与构象隐藏态；频谱定义采用单边角频率规范。\par
\noindent\textcolor{red!65!black}{[跨辑去重未完成 · 草稿]} \par
\medskip\hrule\medskip
\Needspace{7\baselineskip}
\subsubsection*{UC-ADV11-P125 · JMAK成核生长转变的Avrami指数、Arrhenius温移与瞬时转化速率}
\noindent\textbf{题目}\par
某假想晶体相变在300–320 K服从等温Johnson–Mehl–Avrami–Kolmogorov动力学X(t,T)=1−exp\{−[k(T)t]\textasciicircum{}n\}，Avrami指数n=3.000，温度依赖k(T)=k₀exp[−Ea/(RT)]，Ea=60.00 kJ/mol，R=8.314462618 J/(mol K)。在300.0 K测得半转变时间t₅₀=100.0 s，忽略诱导期、晶粒碰撞导致的n变化和潜热反馈。求k300、k320、320 K的t₅₀和t₉₀、320 K半转变时的瞬时dX/dt、两个温度半转变时间之比及k₀；检查X的物理界限与转变速率非负，并解释该模型不能直接推断微观成核机制。\par
\medskip\noindent\textbf{分步解答}
\begin{enumerate}
\item JMAK式X=1−exp[−(kt)\textasciicircum{}n]，0≤X<1，k单位s⁻¹、n为无量纲指数。
\item 半转变条件X=0.5给(kt₅₀)\textasciicircum{}n=ln2，故k300=(ln2)\textasciicircum{}(1/3)/100=0.00884997044501 s⁻¹。
\item Arrhenius温移ln(k320/k300)=(Ea/R)(1/300−1/320)，因此k320=0.0397980755270 s⁻¹，升温加速相变。
\item 320 K半转变时间t₅₀,320=(ln2)\textasciicircum{}(1/3)/k320=22.2371818934 s。
\item 90\%转变条件−ln(1−0.9)=ln10，t₉₀,320=(ln10)\textasciicircum{}(1/3)/k320=33.1800083539 s，必须大于t₅₀。
\item 对JMAK式微分，dX/dt=n k\textasciicircum{}n t\textasciicircum{}(n−1)exp[−(kt)\textasciicircum{}n]，在320 K半转变时速率=0.0467559592679 s⁻¹。
\item 两温度半转变时间之比t₅₀,300/t₅₀,320=k320/k300=4.49697270452。
\item 由k300=k₀exp[−Ea/(RT300)]反推k₀=247551660.992 s⁻¹；其值对Ea假设敏感。
\item 转变速率在t≥0非负且X不会超过1；Avrami指数3不唯一对应某一种晶体生长维度，真实非等温、异质成核和扩散控制需另建模型。
\end{enumerate}
\noindent\textbf{答案摘要}\par
k300 = 0.00884997045 s\textasciicircum{}-1；k320 = 0.0397980755 s\textasciicircum{}-1；t50\_320 = 22.2371819 s；t90\_320 = 33.1800084 s；rate\_at\_half = 0.0467559593 s\textasciicircum{}-1；time\_ratio = 4.4969727 ratio；k0 = 247551661 s\textasciicircum{}-1\par
\noindent\textbf{模型前提：}
等温JMAK、固定Avrami指数、Arrhenius速率与无潜热反馈；不能从n直接断定真实成核几何。\par
\noindent\textcolor{red!65!black}{[跨辑去重未完成 · 草稿]} \par
\medskip\hrule\medskip
\end{document}
